% TO DO ogonek

% \documentclass[a4paper]{article}
% %\usepackage[utf8]{inputenc}
% 
% \begin{document}
% 
% Hi Sadie,
% 
% Did you get this?
% 
% Hope your mood has improved x.
% 
% ---------- Forwarded message ---------
% From: Irena Peterson <irenapetersonpavliuk@gmail.com>
% Date: Sat, Feb 7, 2026 at 10:59 AM
% Subject: THE BORDER GAP (A fragment of my scattered novel)
% To: Sadie <sadie@sadie.org.uk>
% 
% 
% Dear Sadie,
% 
% How are you?
% 
% I'm sending you a very quick, rough fragment of my novel that I typed
% in one breath and then structured using the Bach-inspired
% compositional technique you shared earlier.
% 
% It's part of a larger work that is not yet assembled into a coherent
% whole and is based on real events from my family history. I felt an
% impulse to shape this particular piece and share it with you. I think
% it might be interesting for you as an English reader.
% 
% Don't be surprised if the opening motivation feels a little
% puzzling. In the novel, a new large-scale war in Europe has already
% begun, and one of the key symbolic geographical points is the Suwałki
% corridor in northern Poland, which is occupied in the fictional
% present of the book. I'm still developing this idea and its narrative
% logic.
% 
% The text is included below directly in the email and I'm genuinely
% curious, Would love to hear your thoughts, comments, anything at all!
% 
% Wishing you a beautiful weekend,
% 
% and always glad to be in dialogue with you.
% 
\preludeandfugue{The Border Gap}

\contributor{Irena Peterson Pavliuk}

% in D minor
\prelude

\paragraph{\bf A Moral Polyphony About Luck.}\strut

\begin{quotation}
  
Dear Diary,

\medskip

It's evening. I open the map trying to trace my grandfather's routes in WWII\ldots 

The thought that he was in Poland with the Soviet troops in 1940
suddenly came over me. What was he doing there, in this place that is
all over the news right now? in Suwa{\l}ki\ldots

The road, the only route connecting the Baltic countries with Poland,
has now been blocked there. (Yes, there are actually some forest
roads, but trust me, that's been taken care of too) I think of a small
patch of land between two territories, where enemy armies stand and
shoot at anyone who tries to pass there; I look at it on the
map. Where exactly was my grandfather, who came to the so-called ``new
border'' in 1940 immediately after the mass execution of Polish
prisoners of war in Katyn? The city of Brest figures most prominently
in his memories; it is now a Belarusian city located on the
Polish--Belorussian border\ldots\ I open his memoir, written by my uncle
after my grandfather's death (I have it reprinted on my laptop) and I
begin to read carefully. I must find out the truth. The thought that
my beloved grandfather could have done something stresses me
beyond belief. I know it is impossible, but I will read everything
just to calm myself down.

\end{quotation}

Between borders, between narratives, between names of places and names of guilt,
I look for a human being, a young man who walked through history without knowing
what history would later make of him.

``After graduation in August 1940, Lieutenant P---k was sent to the
Belarusian Military District in the 121st separate anti-tank artillery
division as the commander of a training platoon in which junior
commanders of 45--57~mm anti-tank guns were trained. The division was
located on the territory of Poland along the new border with Germany
along the Western Bug river in the area of the city of Białystok. It
was the famous Suwałki Gap (also known as the Suwałki corridor), from
which the Red Army secretly planned to advance to the west.''

What does it mean ``planned to advance to the west''? The only thing
that calms me down a little in this story is that the Katyn execution
of Polish prisoners of war happened in the spring, April and May 1940,
and my grandfather arrived in Poland only in August. He was barely
twenty years old, I doubt that he fully understood what was
happening. He must have naively believed in the power and greatness of
the Soviet Union and the narrative of liberation.

\fugue

\paragraph{Verse I, Motif I: Barefoot / A Single Pair of Boots.}

So, let's talk about my grandfather. The memoir written by my uncle
begins with the fact that my grandfather grazed pigs in the river
valley as a child. It is strange, but in those years pigs were sent to
the pasture like cows. My little grandfather walked barefoot until the
first snow, and when he went to study at school (the only one of all
his brothers and sisters), he wore the only pair of boots the family
had. He studied excellently. A handwritten letter of commendation
preserved in his archive testifies to his success. It says: ``The
triangle of the six-year Konyshchiv school presents you with this
letter in honor of your energetic dedication to education and public
work. You were in the first rank of the strikers for science. We are
sure that you will continue to set examples for all students''. Dated
December 4, 1933\=the year of the Great Famine.

After that, my grandfather had few options if he wanted to continue
his education. He went to a town 15~kilometers from the village on
foot through forests and fields to study at the gymnasium. He lived in
a dormitory; for this, his family had to tighten their belts, his
mother collected food for her son from all over the village and it was
barely enough for five days; for the whole week it was a loaf of
bread, some fruit and vegetables. Maybe that's why he tried so hard
to provide me with everything I needed when I went to school because
he didn't have anything himself as a child. He would always buy me
these lovely checkered notebooks with pretty birds on the covers.

He very soon understood that the family could no longer support him
and chose to enter the only and very prestigious place available to
him at that time\=a military school, where boys were fed with state
money, because the Soviet Empire always had money for war.

\paragraph{Verse II: War Machinery.}\strut

Moral Motif: Heroism

Counter-Theme: Naivety

My grandfather was a lucky man. He could have been involved in the
Russian-Finnish war as early as 1939, because boys from his school
were taken to the front, but he was lucky to avoid this, as well as
the horrors of the execution of Poles.

So, the Suwałki Gap was\ldots ``from which the Red Army secretly planned
to advance west''. I read further about how they lived in the forest in
dugouts. All recruits of the training platoon came from Central
Asia. They hardly knew the Russian language; my grandfather also had
only basic operation of it which was enough to give orders. ``At the
beginning of~1941, Lieutenant P---k was transferred to the 36th~howitzer
regiment of the 36th~tank division to the position of platoon
commander of the headquarters battery.''

Now it's getting more interesting\ldots ``Later he moved to a rented
apartment from the forest to the city. He fondly remembered the local
Poles, Jewish hairdressers and photographers, their customs, the old
wine, which was not intoxicating, but made his legs sore.'' I so want
to believe that my grandfather was sincere. I really don't understand
all the political vicissitudes, but we have some Polish roots and we
are closer to the Poles in terms of identity, we understand their
language. My grandfather, a young 20-year-old, describes how he was
well received by the Poles despite the fact that he was a lieutenant
in the Soviet Army. I won't question it, I will just leave it here as
it is in the memoir.

After that, the memoir describes the war under the name ``Great
Patriotic War'', as in Soviet textbooks, and it begins, of course, with
the fact that ``on June 22, at dawn, the Germans began intensive
artillery shelling and bombardment from airplanes.'' My grandfather
defended the country he believed in. ``The battery fired at the
Germans, defending the cities of Nesvizh (in modern Belarus), Kovel
(in modern Ukraine), Slutsk, Cherven, the crossing over the Berezina
River\ldots\ For the city of Mogilev in Belarus, battles were fought
for 18 days. In the Smolensk region, the division was encircled, from
which only some units came out fighting.'' This is when grandfather was
appointed to the post of senior officer of the battery. ``The
incessant departure and endless encirclements/entourages (nine in
total) continued.'' After that, I trace on the map the same terrible
road from Mogilev to Moscow, where I see the name Katyn and get chills
over my body\ldots\ My grandfather was not aware of this; he heroically
defended cities of what he considered his country with terribly
Russian names: Yelnya, Yartseve, Vyazma, Tumanov, Gzhatsk,
Mozhaisk\ldots\ all the way to Moscow, and Moscow (God help me
understand this!)~too. He recalled that the most terrible thing was
the mortar shelling, from which there was nowhere to hide. During one
air raid, he hid under the wide steel bed of his 152-millimeter
howitzer\ldots\ But that was exactly what they were bombing. He
remained unharmed, but later, during subsequent raids, he ran away
from guns as far as possible.

In one of the towns with a name I can't even pronounce, my grandfather
was contused and wounded. For his sacrifice and dedication he received
awards; the entire division and the artillery regiment were assigned
the high rank of Guards.

In the fall of~1941, my grandfather's military unit found itself in
its last ``cauldron'' or ``kettle'' (a type of encirclement/entourage in
Soviet military doctrine) near Vyazma. They fought back while there
was ammunition. And then it snowed\ldots\ Early in the morning, they
heard dogs barking. It turned out that the Germans were pulling out
those who were still alive from the trenches. All the prisoners were
taken by convoy to the town of Dorogobuzh in Smolensk region. A
temporary camp was created on the bank of the Dnieper in an abandoned
monastery. There, in the evening, grandfather recognized a mayor he
knew in a prisoner who had spent three weeks there and looked like a
skeleton. The mayor gave him some advice: ``Look, no one gives us food
here. If you want to live, run away while you still can.'' At night,
grandfather and his friend jumped over a gap in the stone wall and
rushed into the sedge growing wild over the Dnieper's bank. The
Germans fired several automatic rounds, but the bullets missed. The
fugitives walked along the bank out of the town and began to wander
through the forests. They ate potatoes dug up in gardens and the
remains of wild apples. At times, they managed to spend a night in
some village. One day grandpa's friend went missing. After that, he
decided to make his way back to Ukraine\ldots

He walked about a thousand kilometers to the south in winter. Since
childhood, he knew nature and how to interact with its forces. He got
food for himself. His mother, my great-grandmother, saved all eight of
her children and several grandchildren during the Great Famine. She
made pancakes from chenopodium, commonly known as the goosefoot, which
occurs almost anywhere in the world, and other weeds.

Everywhere he met people who helped him to hide. The memoir mentions
how, not far from Kyiv, he asked his aunt in the guise of her husband
to pull her cart across the bridge over the Dnieper. He must have
looked like a wild man living in the woods. In Kyiv, which was
occupied by Nazis, he dared to take a risk and went to a barber's.
(It's very strange, because where did he get the money?) The barber
betrayed him: he noticed the officer's boots and called the German
police.

This is what happened next\ldots\ ``The police caught him and locked
him up with other fugitives in a barn. At night, he instigated an
escape. They climbed into the attic, dismantled the thatched roof and
ran away.'' Grandfather then took another risk: he asked a steam
locomotive driver for a ride and this way he almost reached Zhmerynka,
a town in his region. Of course, before getting to town he cautiously
jumped out in the middle of the field and walked through the woods to
his native village.

I wonder, reading this, how he was not afraid of the tribunal or the
fact that his native village was occupied by the Nazis? However, he
was lucky again, there were no Germans in the village, and the
policemen were familiar villagers. It was explained to me when I asked
my uncle that there were many fugitives like my grandfather; the
Germans often allowed prisoners to return home.

This is when the strangest, and to me most incomprehensible, risk of all
took place: my grandfather went back to Brest, Belarus, to pick up his
belongings, which he left in the apartment in 1941. He was lucky: the
landlady did not say a word to the policeman who lived nearby, but
advised him to run away with the suitcase as soon as possible; and so
he did. ``Apparently, he was remembered and loved'', says the memoir.

\paragraph{Verse III, Motif III: Grandmother.}

Grandma Hanna grew up in the occupied village and met my grandpa when
he arrived back after all this adventure; they were around 19--22 years
old and just like everyone of their age attended dance parties. Young
Hanna was only sixteen when Germans arrived in the village but
strangely enough she did not bear any negative memories. I don't
remember exactly from her story, if a young German soldier stayed in
her family house (rented a room or something) or if he only came to
eat with them, but from her I never heard anything about violence in
the village. Still, it seems bizarre. In the memoir it says that my
grandmother remembered the arrival of the partisans, especially a
young Czech named Vyacheslav; her parents helped the partisans in
whatever way they could. It seems that in honour of this Czech, my
grandmother named her son, my uncle.

The main fear of rural youth was that they would be taken to work in
Germany. Young people were transported en masse on big trucks into the
unknown. No one had returned from there yet, so she had no way of
knowing where exactly they went. To avoid this, she entered an
agricultural technical school, since students were not taken to
Germany. This is written twice in the memoir: both about the
grandfather and the grandmother, from which I can conclude that either
they both studied at this technical school, or only one of them, but
who\=it will be difficult to find out.

My grandma's brother, three years younger than her, could not escape
the fate that befell most of the youth in the village. My great uncle
Vlodko repeatedly jumped off the train when they tried to take him to
work in Germany. ``Finally, the Germans got fed up with this,'' my
uncle writes in the memoir, ``and he was sent to the Dachau
concentration camp. He (rather reluctantly) shared how he was
subjected to medical experiments on survival in different temperature
conditions. After the liberation, he served in Poland and in
Germany. He described an interesting incident that occurred to him in
Poland, when in the middle of a small town, he heard someone shout out
his family name. He turned his head and responded; however, it was
addressed not to him, but a local Pole.'' The memoir concludes that
this case is a confirmation of the Polish origin of our family;
``possibly noble roots'', which looks rather challenging against the
background of the whole storyline.

\paragraph{Verse IV: Escape.}

But how did my grandfather avoid being tried for escaping from
captivity?  When the Red Army came to the village in~1944 and battles
were fought under our Raven forest and on the field roads, my
grandfather immediately went to the military and told his name. He
also mentioned that in 1941 he was recommended to the order, but the
documents were lost in the entourages. For two weeks, he was checked
by SMERSH counter-intelligence in a camp in a nearby town. Then he was
called to the officer, who to everyone's amazement stated that
grandfather fought honestly in~1941, and there are no complaints
against him. The officer shook his hand, wished him luck and sent my
grandfather to fight in the 2nd~Ukrainian Front as a platoon
commander.

This is an incredible story, followed by victorious battles proceeding
to the west. He fought in the Carpathians on many ridges; saw horrible
things in abandoned houses such as once the body of a young girl sawn
in two. In Transcarpathia, they liberated Mukacheve, Chop; crossed
over Latorica river and entered Slovakia; then the division was
deployed to the High Beskydy Mountains (the Beskids). In Poland, he
took part in a difficult battle for the liberation of Rudzica railway
station, big cities of Katowice and Częstochowa; forded the Vistula
(where he almost drowned!), then the Oder. In Czechoslovakia they were
greeted with joy: he remembered locals showing them passages in
minefields, and the most beautiful memory was when in the town of
Mladá Boleslav, old women came out with pots of holy water and
sprinkled it over the soldiers with blessings. This is very much at
odds with the modern narrative that nobody wanted to see Soviet troops
in Eastern Europe, but Grandpa said that was indeed the case in
Hungary. In one of the Czech villages, he remembered, local dwellers
approached them with heavy bags, and to the amazement of the hungry
soldiers started taking out sausages and feeding them.

\paragraph{Motif III, Refrain: Grandmother.}

After the liberation of the region by the Soviet Army in 1943--1944,
grandma Hanna went to work in the district committee of the Komsomol,
from where she was sent to study at the Drohobych Pedagogical
Institute in the Lviv region. It is still unclear how to talk about it
and whether this area was liberated from the Nazis by the Red Army,
which followed by joining the region with the rest of Ukraine under
the Soviet empire, or whether it brought only more suffering to its
people as a result of occupation. What is not to be disputed is that
it was a time of great poverty and starvation. Grandma told us how
with her girlfriends, fellow students, they went to the abandoned
fields and private gardens to dig potatoes, left behind by their
deported owners during the infamous ``Vistula'' operation. It was then
when native Poles (and the Ukrainians who had been suspected to
support the Organisation of the Ukrainian nationalists) were deported
to present-day Poland, and Ukrainians who lived to the west of what
the Soviet Union decided to call the border were forcibly transported
from there. Another grief of thousands of people who lost their homes,
on the conscience of the Soviet Empire\ldots\ Then luckily, my young
grandma found a part-time job as a secretary-typist to improve her
financial situation a little.

After graduating from the institute, she married my
grandfather\=Captain P---k, a respected military officer who,
meanwhile, finished his journey almost in Berlin itself.

When the war ended Grandpa received several orders, medals and even
thanks from Stalin, even though shared with us in secret that officers
who survived encirclements/entourages, and especially those of
Ukrainian origin, were not pampered with awards and new military
ranks. Then part of his grandfather's division was demobilized, and he
and his battery were added to the regiment of a rifle division which
was transferred by train to the city of Akhaltsikhe on the
Georgian-Turkish border. The officers said that Stalin was planning a
revenge with Turkey\ldots\ What exactly does this mean? Is it related
to the deportation of Meskhetian Turks?\ldots\ I read that at that time
a total of 115,500 innocent Meskhetian Turks were deported along with
Kurds, Hemshin and Meskhetian Muslims to Central Asia, accused of
counter-revolutionary activities, treason and other nonsense.

Soon grandpa was allowed a long leave. With numerous German trophies,
he came to his native village, helped his father build a new house,
and married my grandmother. The two of them returned to serve in the
Caucasus, which they remember as the best time of their life. In 1947,
there was another famine in Ukraine. My grandfather and grandmother
lived in the Caucasus well, happily, ate deliciously, remembering that
they received an officer's food ration with smoked fish\ldots\ They
didn't even hear about hunger, because of censorship, in the letters
they received from the family, it was painted over in black.

\paragraph{Coda.}

Is grandpa, my hero, someone else's villain? The memoir did not
respond to that question and never will. Any reader will make
conclusions according to their upbringing, that makes a core of values
and beliefs.

It is me who has to respond, interpret the facts of his life, fold
them into a narrative of liberation, occupation\ldots\ then rewrite it
again\ldots

What would a Polish person say? What would a German say? A Kurd? A
Czech? Another Ukrainian?

He was there, doing what his sense of duty required him to. He was a
human. And a very, very lucky one.

% \end{document}
